\chapter{Console command reference}
\label{app:console}

The console is MeerK40t's command line, and it is not a side door: every button,
menu item and pane action ends up as a console command, and many things the
interface does not offer at all are console commands. This appendix is the lookup
table to keep open while you type. Chapter~\ref{ch:console} is the tutorial that
teaches the console; read that first if you have never used a command line.

Open the \panel{Console} window with the ribbon's \btn{Console} button, or toggle
the pane from \menu{Panes>Console} (panes are described in
Section~\ref{sec:panes}). \key{F12} is a popular key to bind to
\cmd{window open Console} in the Keymap window, but it has no binding by default
(Appendix~\ref{app:shortcuts}). In the
console window, \key{Tab} completes command names, \key{Up} and \key{Down} walk
the history, and \key{Enter} runs the line. Three commands answer almost every
question you will have while typing:

\begin{itemize}[leftmargin=1.4em]
  \item \cmd{help} lists every command, grouped by the context it belongs to;
        \cmd{help <command>} prints that command's arguments and options.
  \item \cmd{find <text>} searches command and argument names, e.g.
        \cmd{find dpi}.
  \item \cmd{load_types} and \cmd{save_types} list every file format this
        installation can read and write.
\end{itemize}

\noindent The tables below follow the grouping of \cmd{help}. Every row is one
complete command line that you can type as printed once you have replaced the
placeholders with your own values.

\section*{How to read the syntax column}

\begin{itemize}[leftmargin=1.4em]
  \item \cmd{command} --- literal text: type it exactly like this.
  \item \cmd{<value>} --- a required value that you supply.
  \item \cmd{[<value>]} --- the same, but optional.
  \item \cmd{[--long/-s <value>]} --- one option with two spellings; use either
        \cmd{-}\cmd{-long} or \cmd{-s}, and put a space before its value. Options may
        appear before or after the plain arguments.
  \item \cmd{a|b} --- choose one of the alternatives.
  \item \cmd{...} --- the preceding pattern repeats (used for point lists and for
        multi-part commands).
\end{itemize}

\noindent The brackets, angle brackets, vertical bars and dots are notation and
are never typed. Curly braces mean something else entirely: in a text element they
mark a wordlist reference, as in \cmd{text "Job {serial}"}.

\section*{Session, help and scripting}

{\small
\setlength{\LTleft}{0pt plus 1fill}
\setlength{\LTright}{0pt plus 1fill}
\begin{longtable}{@{}p{3.1cm}p{6.2cm}p{5.1cm}@{}}
\toprule
\textbf{Command} & \textbf{Syntax} & \textbf{What it does} \\
\midrule
\endfirsthead
\toprule
\textbf{Command} & \textbf{Syntax} & \textbf{What it does} \\
\midrule
\endhead
\bottomrule
\endlastfoot
{\footnotesize\cmd{help} \cmd{?}} &
\cmd{help [<command>] [--input <context>] [--output <context>]} &
List all commands, or the arguments and options of one command. \\
{\footnotesize\cmd{find} \cmd{??}} &
\cmd{find <text>} &
List every command and argument name containing \cmd{<text>}. \\
{\footnotesize\cmd{echo}} &
\cmd{echo <text>} &
Print \cmd{<text>} in the console. Mainly useful in scripts, which otherwise run
silently. \\
{\footnotesize\cmd{timer}} &
\cmd{timer[<name>] <count> <seconds> <command>} &
Run \cmd{<command>} \cmd{<count>} times, \cmd{<seconds>} apart. A named timer can be
cancelled with \cmd{timer<name> off}. Bare \cmd{timer} lists the running timers.
Options: \cmd{-}\cmd{-off/-o}, \cmd{-}\cmd{-gui/-g}, \cmd{-}\cmd{-quiet/-q}. \\
{\footnotesize\cmd{loop}} &
\cmd{loop <first> <last> [--variable/-v <name>] <command>} &
Repeat \cmd{<command>} once for every integer from \cmd{<first>} to
\cmd{<last>}. The loop number is written into the command as a brace reference
(default name \cmd{loop}). On Python 3.12 and newer the substitution fails
silently and the literal braces are passed through, so check one loop before you
rely on it. \\
{\footnotesize\cmd{threaded}} &
\cmd{threaded <command>} &
Run the rest of the line in the background so the console is usable again
immediately. This is the prefix the batch and simulation buttons use. \\
{\footnotesize\cmd{wait}} &
\cmd{wait <seconds>} &
Wait before running the next command. This blocks the console, so it belongs in
scripts, not in things you type interactively. \\
{\footnotesize\cmd{quit} \cmd{shutdown}} &
\cmd{quit} &
Shut the program down. Both spellings exist. \\
{\footnotesize\cmd{restart}} &
\cmd{restart} &
Shut down, exit and start MeerK40t again. \\
{\footnotesize\cmd{context}} &
\cmd{context} &
List the context names. A context is a settings scope with its own variables. \\
{\footnotesize\cmd{service}} &
\cmd{service [<domain>]} &
List the services and providers, or enter a service domain. \\
{\footnotesize\cmd{service start}} &
\cmd{service <domain> start <name> [--label/-l <label>] [--path/-p <path>] [--init/-i]} &
Create and start a service from a registered provider, for example
\cmd{service device start grbl}. \\
{\footnotesize\cmd{service activate}} &
\cmd{service <domain> activate <index>} &
Make the service at \cmd{<index>} the active one of its domain. \\
{\footnotesize\cmd{service destroy}} &
\cmd{service <domain> destroy <index>} &
Destroy the service at \cmd{<index>}. \\
{\footnotesize\cmd{module}} &
\cmd{module [(open|close) <name>] [--path/-p <path>]} &
List the open modules, or open and close one. \\
{\footnotesize\cmd{channel}} &
\cmd{channel open <name> [--force/-f]} &
Watch a kernel channel in the console. There is no channel called
\cmd{console}, and unknown names are refused unless you force them. \\
{\footnotesize\cmd{channel close}} &
\cmd{channel close <name>} &
Stop watching a channel. \\
{\footnotesize\cmd{channel list}} &
\cmd{channel list} &
List the channels and mark the ones being watched. \\
{\footnotesize\cmd{channel print}} &
\cmd{channel print <name> [--close/-c]} &
Print a channel to standard output (starting a terminal run). \\
{\footnotesize\cmd{channel save}} &
\cmd{channel save <name> [--filename/-f <file>]} &
Record a channel to disk, by default a timestamped file in the working
directory. \\
{\footnotesize\cmd{execute}} &
\cmd{execute <file>} &
Run every line of \cmd{<file>} as a console command. Blank lines are skipped and
lines starting with \texttt{\#} are comments. \\
{\footnotesize\cmd{batch}} &
\cmd{batch} &
List the stored batch lines with their indices. \\
{\footnotesize\cmd{batch add}} &
\cmd{batch add <line>} &
Append one command to the stored batch list. \\
{\footnotesize\cmd{batch remove}} &
\cmd{batch remove <index>} &
Delete the batch line at \cmd{<index>} (indices are the ones \cmd{batch} prints). \\
{\footnotesize\cmd{batch run}} &
\cmd{batch run <index>} &
Run one batch line. \\
{\footnotesize\cmd{batch enable}} &
\cmd{batch disable|enable <index>} &
Disable or re-enable a batch line without deleting it. \\
{\footnotesize\cmd{alias}} &
\cmd{alias} &
List the defined aliases and the keystroke bound to each. \\
{\footnotesize\cmd{alias set}} &
\cmd{alias <name> <commands>} &
Define an alias. Several commands are separated by \cmd{;}, and the alias is then
a single word you can type (or bind). \\
{\footnotesize\cmd{alias unset}} &
\cmd{alias <name>} &
Delete that alias. \\
{\footnotesize\cmd{alias default}} &
\cmd{alias default} &
Restore the built-in aliases. Note that this discards your own. \\
{\footnotesize\cmd{bind}} &
\cmd{bind} &
List the key bindings and the commands they run. \\
{\footnotesize\cmd{bind set}} &
\cmd{bind <key> <command>} &
Bind a keystroke to a command. An empty command unbinds the key, and
\cmd{bind default} restores the built-in keymap. \\
\bottomrule
\end{longtable}}

\begin{notebox}[Scripts run without feedback]
\cmd{execute} and the batch list run line by line and never wait for you. Anything
that needs a pause (a move, a laser pulse, a human) works better as a plan stage:
add the matching utility operation to the tree rather than relying on
\cmd{wait}. A file run with \cmd{execute} is also the easiest way to keep a job
recipe under version control.
\end{notebox}

\section*{Settings and configuration}

{\small
\setlength{\LTleft}{0pt plus 1fill}
\setlength{\LTright}{0pt plus 1fill}
\begin{longtable}{@{}p{3.1cm}p{6.2cm}p{5.1cm}@{}}
\toprule
\textbf{Command} & \textbf{Syntax} & \textbf{What it does} \\
\midrule
\endfirsthead
\toprule
\textbf{Command} & \textbf{Syntax} & \textbf{What it does} \\
\midrule
\endhead
\bottomrule
\endlastfoot
{\footnotesize\cmd{set}} &
\cmd{set [--path/-p <context>] [<key> <value>]} &
With no key, list the scalar variables of a context (the kernel root by default).
With a key and a value, change that variable. \\
{\footnotesize\cmd{flush}} &
\cmd{flush} &
Write the persistent settings to disk immediately, instead of waiting for the
program to exit. \\
{\footnotesize\cmd{setting_export}} &
\cmd{setting_export <directory>} &
Write a timestamped copy of the whole configuration into \cmd{<directory>}. \\
{\footnotesize\cmd{setting_import}} &
\cmd{setting_import <file>} &
Read a configuration file written by \cmd{setting_export} back in. \\
\bottomrule
\end{longtable}}

\begin{warnbox}[There is no \texttt{get} command]
There is no top-level \cmd{get} or \cmd{setting} command: to read a value, run
\cmd{set} with nothing after it and look through the listing. \cmd{set} also only
writes keys that already exist on that context, and it fails silently when the key
name is wrong, so list first:

\begin{lstlisting}[style=konsole]
set -p elements
set -p elements classify_new False
set -p elements
flush
\end{lstlisting}

\noindent\texttt{get} and \texttt{set} do exist as wordlist subcommands; see the
materials and wordlist table below. Changes made with \cmd{set} are permanent:
they are stored in the configuration and reloaded on the next start.
\end{warnbox}

\section*{Files and documents}

{\small
\setlength{\LTleft}{0pt plus 1fill}
\setlength{\LTright}{0pt plus 1fill}
\begin{longtable}{@{}p{3.1cm}p{6.2cm}p{5.1cm}@{}}
\toprule
\textbf{Command} & \textbf{Syntax} & \textbf{What it does} \\
\midrule
\endfirsthead
\toprule
\textbf{Command} & \textbf{Syntax} & \textbf{What it does} \\
\midrule
\endhead
\bottomrule
\endlastfoot
{\footnotesize\cmd{load}} &
\cmd{load <file>} &
Load a project or drawing. The reader is chosen from the file, so the format
usually does not have to be named. \\
{\footnotesize\cmd{load_types}} &
\cmd{load_types} &
List every format this installation can load. \\
{\footnotesize\cmd{save}} &
\cmd{save <file> [--version/-v <version>]} &
Save the whole scene to \cmd{<file>}. The version must be one of the versions
\cmd{save_types} prints, and it has to match the file extension: the plain and
default versions write \file{.svg}, the compressed version writes \file{.svgz}. \\
{\footnotesize\cmd{save_types}} &
\cmd{save_types} &
List the formats and versions that can be saved. \\
{\footnotesize\cmd{file_startup}} &
\cmd{file_startup} &
Run the stored file-startup commands, the list you edit in the
\panel{AutoExec} pane. \\
{\footnotesize\cmd{xload}} &
\cmd{xload <file> <x> <y> <width> <height>} &
Load a file and place it into the given rectangle on the bed. \\
\bottomrule
\end{longtable}}

\section*{Creating elements}

{\small
\setlength{\LTleft}{0pt plus 1fill}
\setlength{\LTright}{0pt plus 1fill}
\begin{longtable}{@{}p{3.1cm}p{6.2cm}p{5.1cm}@{}}
\toprule
\textbf{Command} & \textbf{Syntax} & \textbf{What it does} \\
\midrule
\endfirsthead
\toprule
\textbf{Command} & \textbf{Syntax} & \textbf{What it does} \\
\midrule
\endhead
\bottomrule
\endlastfoot
{\footnotesize\cmd{rect}} &
\cmd{rect <x> <y> <width> <height> [--rx/-x <r>] [--ry/-y <r>]} &
Rectangle whose top-left corner is at \cmd{<x> <y>}, optionally with rounded
corners. \\
{\footnotesize\cmd{circle}} &
\cmd{circle <x> <y> <r>} &
Circle centred on \cmd{<x> <y>} with radius \cmd{<r>}. \\
{\footnotesize\cmd{ellipse}} &
\cmd{ellipse <cx> <cy> <rx> <ry>} &
Ellipse centred on \cmd{<cx> <cy>}. \\
{\footnotesize\cmd{arc}} &
\cmd{arc <cx> <cy> <rx> <ry> <start> <end> [--rotation/-r <angle>]} &
Arc of an ellipse from \cmd{<start>} to \cmd{<end>}. Write the angles as angles
with a unit, e.g. \cmd{90deg}. \\
{\footnotesize\cmd{line}} &
\cmd{line <x0> <y0> <x1> <y1>} &
Straight line between two points. \\
{\footnotesize\cmd{polyline}} &
\cmd{polyline <x1> <y1> <x2> <y2> ...} &
Open polyline through the listed points. \\
{\footnotesize\cmd{polygon}} &
\cmd{polygon <x1> <y1> <x2> <y2> ...} &
Closed polygon through the listed points. \\
{\footnotesize\cmd{path}} &
\cmd{path "<svg path data>"} &
Path element built from SVG path syntax. Quote the whole path string. \\
{\footnotesize\cmd{text}} &
\cmd{text "<string>" [--size/-s <pt>]} &
Bitmap text element. \\
{\footnotesize\cmd{linetext}} &
\cmd{linetext <x> <y> "<text>" [--font/-f <file>] [--font_size/-s <size>] [--font_spacing/-g <factor>]} &
Single-line vector text drawn with a Hershey font, so it can be cut rather than
engraved. \\
\bottomrule
\end{longtable}}

\begin{notebox}[Elements appear selected]
Every command in this table puts the new element into the selection, so the next
command can act on it, and it is classified into a matching operation if you have
\texttt{classify new} enabled. There is no console command that creates a raster
image element: images arrive through \cmd{load} or the interface's import, which
is why the image commands below take an existing image as their input.
\end{notebox}

\section*{Selecting, transforming and editing elements}

Commands whose first column contains a bare \cmd{element} or \cmd{operation}
prefix take their input from the console's current data rather than from the
selection. Commands without such a prefix, such as \cmd{scale}, start a new chain
and fall back to the current selection.

{\small
\setlength{\LTleft}{0pt plus 1fill}
\setlength{\LTright}{0pt plus 1fill}
\begin{longtable}{@{}p{3.1cm}p{6.2cm}p{5.1cm}@{}}
\toprule
\textbf{Command} & \textbf{Syntax} & \textbf{What it does} \\
\midrule
\endfirsthead
\toprule
\textbf{Command} & \textbf{Syntax} & \textbf{What it does} \\
\midrule
\endhead
\bottomrule
\endlastfoot
{\footnotesize\cmd{element}} &
\cmd{element} &
The currently selected elements, ready to be passed to the next command. \\
{\footnotesize\cmd{element*}} &
\cmd{element*} &
Every element in the scene. \\
{\footnotesize\cmd{element~}} &
\cmd{element~} &
Every element that is not selected. \\
{\footnotesize\cmd{element0,3,4}} &
\cmd{element<n>,<n>} &
The elements with those indices, as printed by \cmd{element* list}. \\
{\footnotesize\cmd{select}} &
\cmd{<data> select} &
Make the data the selection, e.g. \cmd{element* select} to select everything. \\
{\footnotesize\cmd{select+}} &
\cmd{<data> select+} &
Add the data to the selection. \\
{\footnotesize\cmd{select-}} &
\cmd{<data> select-} &
Remove the data from the selection. \\
{\footnotesize\cmd{select^}} &
\cmd{<data> select^} &
Toggle the data in the selection. \\
{\footnotesize\cmd{range}} &
\cmd{<data> range <first> <last> [--step/-s <n>]} &
Reduce the data to a slice by index before passing it on. \\
{\footnotesize\cmd{id}} &
\cmd{<data> id [<id>]} &
Print, or set, the ids of the data. \\
{\footnotesize\cmd{scale}} &
\cmd{scale <sx> [<sy>] [--px/-x <x>] [--py/-y <y>] [--absolute/-a]} &
Scale the selection, optionally differently in \cmd{y} and from a chosen origin
instead of the centre of the selection. \\
{\footnotesize\cmd{rotate}} &
\cmd{rotate <angle> [--cx/-x <x>] [--cy/-y <y>] [--absolute/-a]} &
Rotate the selection around its centre, or around \cmd{<cx> <cy>}. \\
{\footnotesize\cmd{translate}} &
\cmd{translate <tx> <ty> [--absolute/-a]} &
Move the selection by an offset. \\
{\footnotesize\cmd{position}} &
\cmd{position <x> <y>} &
Move the selection so that it starts at \cmd{<x> <y>}. \\
{\footnotesize\cmd{move_to_laser}} &
\cmd{move_to_laser} &
Move the selection to the current laser position. \\
{\footnotesize\cmd{matrix}} &
\cmd{matrix <sx> <kx> <ky> <sy> <tx> <ty>} &
Set the selection's transformation matrix directly. \\
{\footnotesize\cmd{resize}} &
\cmd{resize <x> <y> <width> <height>} &
Move and scale the selection so that it fills the given rectangle. \\
{\footnotesize\cmd{reset}} &
\cmd{reset} &
Reset the affine transforms of the selection back to identity. \\
{\footnotesize\cmd{area}} &
\cmd{area [--new_area/-n <area>] [--density/-d <n>]} &
Report, or change, the area a raster element covers. \\
{\footnotesize\cmd{copy}} &
\cmd{copy [--dx/-x <dx>] [--dy/-y <dy>] [--copies/-c <n>]} &
Duplicate the selection, optionally offset and repeated. \\
{\footnotesize\cmd{delete}} &
\cmd{<data> delete} &
Delete the data, e.g. \cmd{element delete}. \cmd{delete} on its own is rejected,
because it does not accept the base context. \\
{\footnotesize\cmd{group}} &
\cmd{group <label> [--reparent/-r]} &
Group the selection under a new label. \\
{\footnotesize\cmd{ungroup}} &
\cmd{ungroup} &
Dissolve the selected groups. \\
\bottomrule
\end{longtable}}

\noindent Note the asymmetry: \cmd{scale}, \cmd{rotate}, \cmd{translate},
\cmd{position}, \cmd{resize}, \cmd{copy} and \cmd{area} can be typed on their own
and then work on the selection, while \cmd{delete}, \cmd{select},
\cmd{select+}, \cmd{select-}, \cmd{select^}, \cmd{range} and \cmd{id} must be
given their data, for example \cmd{element* select} or \cmd{element delete}.
\cmd{move} is not an element command at all: it moves the laser head, and lives in
the device table below.

\section*{Operations and the plan pipeline}

{\small
\setlength{\LTleft}{0pt plus 1fill}
\setlength{\LTright}{0pt plus 1fill}
\begin{longtable}{@{}p{3.1cm}p{6.2cm}p{5.1cm}@{}}
\toprule
\textbf{Command} & \textbf{Syntax} & \textbf{What it does} \\
\midrule
\endfirsthead
\toprule
\textbf{Command} & \textbf{Syntax} & \textbf{What it does} \\
\midrule
\endhead
\bottomrule
\endlastfoot
{\footnotesize\cmd{cut} \cmd{engrave} \cmd{raster} \cmd{imageop} \cmd{dots} \cmd{hatch}} &
\cmd{<type> [--speed/-s <mm/s>] [--power/-p <ppi>] [--dpi/-d <n>] [--overscan/-o <length>] [--passes/-x <n>] [--color/-c <colour>] [--default/-D] [--parallel/-P] [--stroke/-K] [--fill/-F]} &
Create an operation of that type and attach the elements it is given to it. With
\cmd{-}\cmd{-parallel} every element gets its own operation instead. \\
{\footnotesize\cmd{operation} \cmd{operation*} \cmd{operation~}} &
\cmd{operation}, \cmd{operation*}, \cmd{operation<n>,<n>} &
The selected, all, unselected or listed operations, ready to be passed on. \\
{\footnotesize\cmd{operation* list}} &
\cmd{operation* list} &
List the operations with their children. \\
{\footnotesize\cmd{operation* empty}} &
\cmd{operation* empty} &
Remove every child from the selected operations, leaving the operations. \\
{\footnotesize\cmd{operation* enable}} &
\cmd{operation* disable}, \cmd{operation* enable} &
Disable or enable operations, so they are left out of the plan. \\
{\footnotesize\cmd{plan}} &
\cmd{plan[<name>] <command>} &
Work on the plan called \cmd{<name>}. Bare \cmd{plan} uses the plan named
\cmd{0}; typing \cmd{planz} renames the default plan to \cmd{z} for this session,
and \cmd{plan2} switches to the plan named \cmd{2}. \\
{\footnotesize\cmd{list-plans}} &
\cmd{list-plans} &
List the plans that exist and the stage each one has reached. \\
{\footnotesize\cmd{plan clear}} &
\cmd{plan clear} &
Empty the plan. \\
{\footnotesize\cmd{plan copy}} &
\cmd{plan copy} &
Copy the artwork of every operation into the plan. \\
{\footnotesize\cmd{plan copy-selected}} &
\cmd{plan copy-selected} &
Copy the artwork of the selected operations only. \\
{\footnotesize\cmd{plan preprocess}} &
\cmd{plan preprocess} &
Prepare the plan for execution: convert the scene into device coordinates. \\
{\footnotesize\cmd{plan validate}} &
\cmd{plan validate} &
Run the plan's validation. \\
{\footnotesize\cmd{plan blob}} &
\cmd{plan blob} &
Convert the plan into the device's own instruction format. \\
{\footnotesize\cmd{plan preopt}} &
\cmd{plan preopt} &
Prepare the plan for optimisation. \\
{\footnotesize\cmd{plan optimize}} &
\cmd{plan optimize} &
Optimise the plan, for example to reduce travel between cuts. \\
{\footnotesize\cmd{plan finish}} &
\cmd{plan finish} &
Mark the plan finished, which is what the simulation expects. \\
{\footnotesize\cmd{plan return}} &
\cmd{plan return} &
Push the plan's operations back into the tree. \\
{\footnotesize\cmd{plan sublist}} &
\cmd{plan sublist [<first>] [<last>]} &
Keep only that slice of the plan, where the last index \cmd{-1} means the end. \\
{\footnotesize\cmd{plan console}} &
\cmd{plan console [--index/-i <n>] <command>} &
Insert a console command into the plan at that position. \\
{\footnotesize\cmd{spool}} &
\cmd{spool} &
Send the plan it is given to the device spooler. Typed on its own it prints the
spooler queue instead. \\
{\footnotesize\cmd{save_job}} &
\cmd{plan <stages> save_job <file>} &
Export the plan in the active driver's own file format instead of burning it.
Only drivers that have such a format offer this. \\
\bottomrule
\end{longtable}}

\begin{description}[leftmargin=1.2em,style=nextline]
  \item[The order matters.] \cmd{copy}, \cmd{preprocess}, \cmd{validate},
        \cmd{blob}, \cmd{preopt}, \cmd{optimize} must run in that order;
        \cmd{spool} refuses a plan that has not been through \cmd{blob}, because
        there is nothing in the device's language yet. Each stage records itself,
        which is what \cmd{list-plans} shows you.
  \item[\cmd{copy} or \cmd{copy-selected}?] \cmd{copy} takes the artwork of every
        operation that produces output. \cmd{copy-selected} takes only the
        selected ones, which is how the interface burns a single operation.
\end{description}

\section*{Device, motion and the spooler}

{\small
\setlength{\LTleft}{0pt plus 1fill}
\setlength{\LTright}{0pt plus 1fill}
\begin{longtable}{@{}p{3.1cm}p{6.2cm}p{5.1cm}@{}}
\toprule
\textbf{Command} & \textbf{Syntax} & \textbf{What it does} \\
\midrule
\endfirsthead
\toprule
\textbf{Command} & \textbf{Syntax} & \textbf{What it does} \\
\midrule
\endhead
\bottomrule
\endlastfoot
{\footnotesize\cmd{device}} &
\cmd{device} &
List the device providers that can be added, with the family each one belongs
to. \\
{\footnotesize\cmd{device add}} &
\cmd{device add <provider> [--label/-l <label>]} &
Create and start a new device from a provider, e.g.
\cmd{device add grbl-generic}. \\
{\footnotesize\cmd{device activate}} &
\cmd{device activate [<name>|<index>]} &
With no argument, list the defined devices with their labels, types and status.
With a label or an index, make that one the active device. \\
{\footnotesize\cmd{device duplicate}} &
\cmd{device duplicate <name>|<index> [<label>]} &
Copy an existing device entry, which is how you keep two configurations of the
same machine. \\
{\footnotesize\cmd{device delete}} &
\cmd{device delete <name>|<index>} &
Delete a device entry. The active device is protected. \\
{\footnotesize\cmd{devinfo}} &
\cmd{devinfo} &
Print the current position in the form \cmd{x,y;nx,ny;}. \\
{\footnotesize\cmd{home}} &
\cmd{home} &
Home the laser, according to the device's own definition of home. \\
{\footnotesize\cmd{physical_home}} &
\cmd{physical_home} &
Home against the physical endstops. \\
{\footnotesize\cmd{lock} \cmd{unlock}} &
\cmd{lock}, \cmd{unlock} &
Lock or unlock the rail, so the head cannot be pushed around by accident. \\
{\footnotesize\cmd{move} \cmd{move_absolute}} &
\cmd{move <x> <y> [--force/-f]} &
Move the head to an absolute position. Without \cmd{-}\cmd{-force} the move is refused
while the spooler is busy (\cmd{Busy Error}). \\
{\footnotesize\cmd{move_relative}} &
\cmd{move_relative <dx> <dy> [--force/-f]} &
Move the head by an offset. Same busy rule. \\
{\footnotesize\cmd{left} \cmd{right} \cmd{up} \cmd{down}} &
\cmd{left <amount>} &
Queue a relative jog, \unit{1}\unit{mm} when no amount is given. The move is sent
a moment later by \cmd{jog}, which is how the arrow keys produce a continuous
nudge while held. \\
{\footnotesize\cmd{jog}} &
\cmd{jog [--force/-f]} &
Send the queued jog move. \cmd{-}\cmd{-force} sends it even when the device would
otherwise confine it. \\
{\footnotesize\cmd{+laser} \cmd{-laser}} &
\cmd{+laser}, \cmd{-laser} &
Turn the laser on or off without moving the head. \\
{\footnotesize\cmd{spool list}} &
\cmd{spool list} &
List the jobs in the spooler queue. \\
{\footnotesize\cmd{spooler clear}} &
\cmd{spooler clear} &
Empty the spooler queue. \\
{\footnotesize\cmd{spooler send}} &
\cmd{spooler send <command>} &
Send a single plan command (\cmd{home}, \cmd{unlock}, \cmd{origin}, \ldots) to the
spooler as a small job. \\
{\footnotesize\cmd{pause} \cmd{resume}} &
\cmd{pause}, \cmd{resume} &
Pause or resume the running job. Only registered by drivers that can hold a job. \\
{\footnotesize\cmd{estop} \cmd{abort}} &
\cmd{estop}, \cmd{abort} &
Abort the running job. Driver dependent. \\
{\footnotesize\cmd{goto}} &
\cmd{goto <x> <y>} &
Send a goto to a galvo controller. Driver dependent. \\
\bottomrule
\end{longtable}}

\begin{warnbox}[Motion needs a device]
The motion and job commands come from the active device's driver, not from the
kernel: a machine without endstops has no \cmd{physical_home} that does anything
useful, and \cmd{pause} does not exist on every driver. If a command in this table
is rejected as ``not a registered command'', check which device is active with
\cmd{device activate} and see Appendix~\ref{app:drivers} for what that driver
supports.
\end{warnbox}

\section*{Materials and wordlists}

{\small
\setlength{\LTleft}{0pt plus 1fill}
\setlength{\LTright}{0pt plus 1fill}
\begin{longtable}{@{}p{3.1cm}p{6.2cm}p{5.1cm}@{}}
\toprule
\textbf{Command} & \textbf{Syntax} & \textbf{What it does} \\
\midrule
\endfirsthead
\toprule
\textbf{Command} & \textbf{Syntax} & \textbf{What it does} \\
\midrule
\endhead
\bottomrule
\endlastfoot
{\footnotesize\cmd{material}} &
\cmd{material} &
Enter the materials context, which holds the sections of your material library. \\
{\footnotesize\cmd{material save}} &
\cmd{material save <name> [--author/-a <author>] [--description/-d <text>]} &
Save the current materials into the settings under \cmd{<name>}. \\
{\footnotesize\cmd{material load}} &
\cmd{material load <name>} &
Load a saved material library. \\
{\footnotesize\cmd{material list}} &
\cmd{material list [<name>]} &
Show the material sections, or one named library. \\
{\footnotesize\cmd{material delete}} &
\cmd{material delete <name>} &
Delete a saved material library. \\
{\footnotesize\cmd{wordlist}} &
\cmd{wordlist} &
Enter the wordlist context. \\
{\footnotesize\cmd{wordlist add}} &
\cmd{wordlist add <key> [<value>]} &
Add a key with a value (empty when none is given). \\
{\footnotesize\cmd{wordlist addcounter}} &
\cmd{wordlist addcounter <key> [<start>]} &
Add a numeric counter that starts at \cmd{<start>}, or 1. \\
{\footnotesize\cmd{wordlist get}} &
\cmd{wordlist get <key> [<index>]} &
Print the value the key has at that index. \\
{\footnotesize\cmd{wordlist set}} &
\cmd{wordlist set <key> <value> [<index>]} &
Write a value at an index. \\
{\footnotesize\cmd{wordlist index}} &
\cmd{wordlist index <key|@ALL> <index|+n>} &
Move an index, where \cmd{+2} means ``two on'' rather than ``set to 2''. \\
{\footnotesize\cmd{wordlist list}} &
\cmd{wordlist list [<key>]} &
List the keys, or the values stored under one key. \\
{\footnotesize\cmd{wordlist advance}} &
\cmd{wordlist advance} &
Move every index on, so the next job gets the next values. Spooling a plan runs
this for you. \\
{\footnotesize\cmd{wordlist restore}} &
\cmd{wordlist restore [<file>]} &
Load a saved wordlist, by default the standard one. \\
{\footnotesize\cmd{wordlist backup}} &
\cmd{wordlist backup [<file>]} &
Save the current wordlist. \\
{\footnotesize\cmd{wordlist load}} &
\cmd{wordlist load <csv-file>} &
Attach a CSV file to the wordlist, which is how you import a list of names or
serial numbers. \\
\bottomrule
\end{longtable}}

\begin{notebox}[How the wordlist reaches the machine]
A text element can contain brace references such as \cmd{{name}} or
\cmd{{serial}}. They are left alone on screen and replaced by the wordlist
values while the plan is built, which is why the preview still shows the braces.
The job that is built is the one that carries the substituted text, so build a
plan and spool it before judging the result.
\end{notebox}

\section*{Images, rendering and tracing}

{\small
\setlength{\LTleft}{0pt plus 1fill}
\setlength{\LTright}{0pt plus 1fill}
\begin{longtable}{@{}p{3.1cm}p{6.2cm}p{5.1cm}@{}}
\toprule
\textbf{Command} & \textbf{Syntax} & \textbf{What it does} \\
\midrule
\endfirsthead
\toprule
\textbf{Command} & \textbf{Syntax} & \textbf{What it does} \\
\midrule
\endhead
\bottomrule
\endlastfoot
{\footnotesize\cmd{image}} &
\cmd{image} &
Enter the image context for the selected images, which is how the raster
commands below are reached. \\
{\footnotesize\cmd{image dewhite}} &
\cmd{image dewhite} &
Remove the white background from the image. \\
{\footnotesize\cmd{image rgba}} &
\cmd{image rgba} &
Convert the image to RGBA. \\
{\footnotesize\cmd{image quantize}} &
\cmd{image quantize [<colours>]} &
Reduce the image to that many colours, 10 by default. \\
{\footnotesize\cmd{image solarize}} &
\cmd{image solarize [<threshold>]} &
Solarise the tones, with a threshold of 250 by default. \\
{\footnotesize\cmd{image save}} &
\cmd{image save <file> [--processed/-p]} &
Write the image to disk, optionally the processed result rather than the
original. \\
{\footnotesize\cmd{render}} &
\cmd{render [--dpi/-d <dpi>]} &
Rasterise the selected elements into an image element, 500\unit{dpi} by
default. \\
{\footnotesize\cmd{render_split}} &
\cmd{render_split <columns> <rows> <dpi> [--order/-o <mode>]} &
Rasterise and cut the result into a grid of image parts, for tiling a large
engraving. \\
{\footnotesize\cmd{render_keyhole}} &
\cmd{render_keyhole <dpi> [--order/-o <mode>] [--invert/-i] [--outline/-b]} &
Rasterise the selection as a keyhole mask, with optional inverted mask and
outline. \\
{\footnotesize\cmd{vectorize}} &
\cmd{vectorize [--dpi/-d <dpi>] [--turnpolicy/-z <policy>] [--turdsize/-t <n>] [--alphamax/-a <f>] [--opticurve/-n] [--opttolerance/-O <f>] [--color/-C <colour>] [--invert/-i] [--blacklevel/-k <f>]} &
Convert the selected elements into paths by rasterising them and tracing the
result. \\
{\footnotesize\cmd{potrace}} &
\cmd{image potrace [<options>]} &
Trace the selected image into paths with Potrace. It takes the image options
above. \\
{\footnotesize\cmd{vectrace}} &
\cmd{image vectrace} &
Trace the selected image into paths by following the boundaries between black
and white. \\
{\footnotesize\cmd{vtracer}} &
\cmd{vtracer} &
Trace the selection into paths with VTracer, which also handles colour images.
It needs the optional \file{vtracer} package installed. \\
{\footnotesize\cmd{trace}} &
\cmd{trace <method> <resolution> [--start/-s <0|1|2>] [--force/-f]} &
Work out the outline of the selection and send that outline to the device as a
job, so you can see where the job would run. Methods: \cmd{quick}, \cmd{hull},
\cmd{segment}, \cmd{circle}. \\
{\footnotesize\cmd{tracegen}} &
\cmd{tracegen <method> <resolution>} &
The same outline, but created as a new element instead of being burned. \\
\bottomrule
\end{longtable}}

\noindent \cmd{trace} and \cmd{tracegen} default to the selection and can be
redirected with a chain. Elements that are not in an enabled operation are left
out of the outline unless you pass \cmd{-}\cmd{-force}, and the console tells you how
many were skipped.

\section*{Windows, panes and dialogs}

These commands need the graphical interface; in a console-only run
(\cmd{-}\cmd{-no-gui}) they are not registered at all.

{\small
\setlength{\LTleft}{0pt plus 1fill}
\setlength{\LTright}{0pt plus 1fill}
\begin{longtable}{@{}p{3.1cm}p{6.2cm}p{5.1cm}@{}}
\toprule
\textbf{Command} & \textbf{Syntax} & \textbf{What it does} \\
\midrule
\endfirsthead
\toprule
\textbf{Command} & \textbf{Syntax} & \textbf{What it does} \\
\midrule
\endhead
\bottomrule
\endlastfoot
{\footnotesize\cmd{window}} &
\cmd{window [--path/-p <context>]} &
List the open and registered windows in that context. \\
{\footnotesize\cmd{window list}} &
\cmd{window list} &
List the window names that can be opened. \\
{\footnotesize\cmd{window displays}} &
\cmd{window displays} &
Show the displays and where each open window currently sits. \\
{\footnotesize\cmd{window open} \cmd{window toggle}} &
\cmd{window open <window> [--multi/-m <n>]} &
Open a window, or close it again if it is open. \cmd{-}\cmd{-multi/-m} with an index
opens another copy. Typical names: \cmd{Console}, \cmd{AutoExec},
\cmd{Preferences}, \cmd{JobSpooler}, \cmd{Simulation}, \cmd{DeviceManager}. \\
{\footnotesize\cmd{window close}} &
\cmd{window close <window>} &
Close the window. \\
{\footnotesize\cmd{window reset}} &
\cmd{window reset <window>} &
Forget the stored position and size of one window, or of all of them with
\cmd{*}. \\
{\footnotesize\cmd{page}} &
\cmd{page [<page>]} &
Switch the ribbon bar to a page, or report which page is active. \\
{\footnotesize\cmd{pane}} &
\cmd{pane} &
Enter the pane context. \\
{\footnotesize\cmd{pane show} \cmd{pane hide} \cmd{pane toggle}} &
\cmd{pane show <pane>}, \cmd{pane hide <pane>}, \cmd{pane toggle <pane>} &
Show, hide, or show-or-hide a docked pane, e.g. \cmd{pane toggle spooler}. \\
{\footnotesize\cmd{pane float} \cmd{pane dock}} &
\cmd{pane float [--always/-a] <pane>}, \cmd{pane dock <pane>} &
Float a pane as a free window, or put it back in the frame. \\
{\footnotesize\cmd{pane save} \cmd{pane load}} &
\cmd{pane save <name>}, \cmd{pane load <name>} &
Store the current pane layout under a name, or restore a stored one. \\
{\footnotesize\cmd{pane reset}} &
\cmd{pane reset} &
Restore the default layout. \\
{\footnotesize\cmd{pane lock} \cmd{pane unlock}} &
\cmd{pane lock}, \cmd{pane unlock} &
Lock the panes so they cannot be moved, or unlock them. \\
{\footnotesize\cmd{pane toggleui}} &
\cmd{pane toggleui} &
Hide every pane except the scene, or bring them all back. \\
{\footnotesize\cmd{dialog_path}} &
\cmd{dialog_path} &
Open a dialog to type SVG path data and add the path to the scene. \\
{\footnotesize\cmd{dialog_transform}} &
\cmd{dialog_transform} &
Open a dialog to type an SVG transform such as \cmd{scale(1.5)} and apply it to
the elements. \\
{\footnotesize\cmd{dialog_fill} \cmd{dialog_stroke}} &
\cmd{dialog_fill}, \cmd{dialog_stroke} &
Open the colour picker for the fill or the stroke of the selection. \\
\bottomrule
\end{longtable}}

\noindent The four dialog commands are hidden from \cmd{help}, because they exist
only to put a keyboard shortcut on an interface dialog. They are listed here so
that a macro of yours can use them.

\section*{Remote control and server commands}

These start servers that let other software drive this copy of MeerK40t, or let
another computer forward data to the laser. \cmd{-}\cmd{-quit/-q} stops the server
again, and \cmd{-}\cmd{-verbose/-v} mirrors the server's own channel into the console.

{\small
\setlength{\LTleft}{0pt plus 1fill}
\setlength{\LTright}{0pt plus 1fill}
\begin{longtable}{@{}p{3.1cm}p{6.2cm}p{5.1cm}@{}}
\toprule
\textbf{Command} & \textbf{Syntax} & \textbf{What it does} \\
\midrule
\endfirsthead
\toprule
\textbf{Command} & \textbf{Syntax} & \textbf{What it does} \\
\midrule
\endhead
\bottomrule
\endlastfoot
{\footnotesize\cmd{grblcontrol}} &
\cmd{grblcontrol [--port/-p <port>] [--verbose/-v] [--quit/-q]} &
Listen on a TCP port (23 by default) and answer as a GRBL controller, so GRBL
software can send jobs through MeerK40t to this machine. \\
{\footnotesize\cmd{ruidacontrol}} &
\cmd{ruidacontrol [--verbose/-v] [--quit/-q] [--jogless/-j] [--man_in_the_middle/-m <host>] [--bridge/-b]} &
The same idea for Ruida software. \cmd{-}\cmd{-man_in_the_middle} forwards the
traffic to a real laser instead of emulating it, and \cmd{-}\cmd{-bridge} uses the
bridge protocol. \\
{\footnotesize\cmd{lhyserver}} &
\cmd{lhyserver [--port/-p <port>] [--verbose/-v] [--watch/-w] [--quit/-q]} &
Serve the lihuiyu (K40) controller over TCP, so a small computer next to the
laser can forward its USB traffic. \cmd{-}\cmd{-watch} also logs the data. \\
{\footnotesize\cmd{webserver}} &
\cmd{webserver [--port/-p <port>] [--silent/-s] [--quit/-q]} &
Start the built-in web server, port 2080 by default. \\
\bottomrule
\end{longtable}}

\section*{Worked command chains}

The console runs one command per line, but a line may hold several commands: a
\cmd{|} between them splits the line, unless the bar sits inside a quoted string.
Commands also chain without any separator at all, because each command passes its
result to the next one. A command that cannot create its own input has to follow a
command that produces it; when a command returns nothing, chaining stops and the
console says so.

\begin{lstlisting}[style=konsole]
rect 0 0 20mm 10mm engrave -s 15 plan copy-selected preprocess validate blob preopt optimize spool
\end{lstlisting}

\noindent Read from the left: \cmd{rect} creates the element and selects it;
\cmd{engrave -s 15} creates an engrave operation at \unit{15}\unit{mm/s} and
attaches the rectangle to it; \cmd{plan} hands that operation to the plan named
\cmd{0}; \cmd{copy-selected} copies only the selected operations into the plan;
\cmd{preprocess} converts the scene into device coordinates; \cmd{validate},
\cmd{blob}, \cmd{preopt} and \cmd{optimize} prepare the device format and reduce
travel; and \cmd{spool} queues the result on the device. Nothing moves until you
press the machine's start button or run \cmd{startjob}.

The same job written one stage per line is easier to read and easier to paste into
the \panel{AutoExec} pane, because each line restates the plan:

\begin{lstlisting}[style=konsole]
rect 0 0 20mm 10mm | engrave -s 15
plan copy-selected
plan preprocess
plan validate
plan blob
plan preopt
plan optimize
plan spool
\end{lstlisting}

\noindent To write the job to a file instead of burning it, replace the last stage
with \cmd{save_job <file>}:

\begin{lstlisting}[style=konsole]
rect 0 0 20mm 10mm engrave plan copy preprocess validate blob preopt optimize save_job job.egv
\end{lstlisting}

\noindent\cmd{save_job} writes whatever the active driver sends to the machine, so
the extension depends on the driver. Use \cmd{plan copy} here rather than
\cmd{copy-selected} only when you want the whole project, not one operation.

The two house aliases do the whole pipeline for you, and \cmd{alias} prints them
along with everything else you have defined:

\begin{lstlisting}[style=konsole]
burn
simulate
alias
\end{lstlisting}

\noindent \cmd{burn} is \cmd{planz clear copy preprocess validate blob preopt
optimize spool}, and \cmd{simulate} is the same chain with \cmd{finish} and the
simulation window instead of \cmd{spool}. Both switch the default plan to \cmd{z}
on the way, so a \cmd{plan} command you type afterwards refers to that plan.

Loading, changing and saving a document is three lines, and the save format is
chosen by the extension:

\begin{lstlisting}[style=konsole]
load drawing.svg
element* select
scale 0.5
save smaller.svg
\end{lstlisting}

\noindent A saved copy of a scene can be compressed by asking for the matching
version, \cmd{save smaller.svgz -v compressed}; the version and the extension must
agree, which is what \cmd{save_types} lists.

Adding a machine is the same three commands every time: list the providers, add
one, and make it active. \cmd{device activate} on its own lists the devices you
already have, with the index the other forms accept:

\begin{lstlisting}[style=konsole]
device
device add grbl-generic
device activate
device activate Grbl
\end{lstlisting}

\noindent Personalising text is a wordlist plus a brace reference. The
substitution happens while the plan is built, so build and spool the job to see
it:

\begin{lstlisting}[style=konsole]
wordlist add name Alice
wordlist addcounter serial 100
wordlist list
text "Hello {name}"
\end{lstlisting}

\noindent Between jobs, \cmd{wordlist advance} steps every index on, and spooling
does that for you. To keep several jobs apart instead, put the whole recipe in a
file and run it with \cmd{execute}, or store it in the batch list with
\cmd{batch add <line>}.

\section*{Gotchas}

\begin{itemize}[leftmargin=1.4em]
  \item \textbf{Always write a unit.} A bare number is not a millimetre and not a
        pixel: for lengths it is a raw internal unit of one 65535th of an inch, so
        \cmd{rect 0 0 10 10} produces a rectangle you cannot see, and for angles
        it is radians, so \cmd{rotate 45} turns the element seven times over
        rather than 45 degrees. Write \cmd{rect 0 0 10mm 10mm} and
        \cmd{rotate 45deg}. \cmd{cm}, \cmd{mm}, \cmd{in}, \cmd{mil}, \cmd{px},
        \cmd{pt} and \texttt{\%} are all accepted; a percentage is a fraction of the
        bed, the width for \cmd{x} and \cmd{width} and the height for \cmd{y} and
        \cmd{height}. In the console of the graphical interface, \cmd{unit 2cm}
        translates one value into every other unit.
  \item \textbf{Elements and operations need selecting.} Transform and
        element commands work on the selection when they start a chain, so
        \cmd{scale 0.5} with nothing selected reports ``No selected elements'' and
        does nothing. The pipeline is no different: \cmd{plan copy-selected}
        copies the selected operations, and if none is selected the plan is
        quietly empty. Run the selection commands first
        (\cmd{element* select}, \cmd{operation* select}).
  \item \textbf{Some commands cannot start a line.} \cmd{delete}, \cmd{select},
        \cmd{range}, \cmd{id} and the \cmd{operation* list} family only accept
        elements, operations or a plan as input. Typed first they are refused
        (``not a registered command in this context: Base''); put the input in
        front of them: \cmd{element* select}, \cmd{element delete},
        \cmd{operation* list}.
  \item \textbf{Quote anything with spaces.} Only double quotes group a value;
        single quotes are not special, and there are no escape sequences. A comma
        also separates arguments, so quote text that contains one:
        \cmd{text "Hello world"}.
  \item \textbf{Write options with a space.} \cmd{-s 15} is the speed option with
        value 15. \cmd{-s15} is read as three separate options and will not do
        what you meant.
  \item \textbf{Case.} Command and option names are matched in lower case, so
        \cmd{RECT} works, but the values you pass keep their case, which matters
        for text and file names.
  \item \textbf{\cmd{set} is scoped and silent.} \cmd{set -p <context> <key>
        <value>} changes a variable that must already exist on that context; a
        wrong key name does nothing at all and prints nothing. List with
        \cmd{set -p <context>} first, and use \cmd{flush} to write settings to
        disk. The change is permanent until you change it back.
  \item \textbf{Plan names stick.} \cmd{planz} is not a typo: the letters after
        \cmd{plan} become the name of the default plan for the rest of the
        session, and bare \cmd{plan} then means ``the plan named z''. The plan
        that a fresh session uses is named \cmd{0}.
  \item \textbf{The pipeline is ordered.} \cmd{spool} refuses a plan that has not
        been through \cmd{blob}, and \cmd{optimize} does nothing useful before
        \cmd{blob} has run. When a chain half-works, \cmd{list-plans} shows how
        far the plan got.
  \item \textbf{Drivers differ.} Anything that moves the head or changes the state
        of the machine belongs to the active device's driver. \cmd{pause},
        \cmd{estop}, \cmd{status} and \cmd{goto} exist only where a driver
        implements them, so a command that works on one machine may be rejected
        on another.
\end{itemize}
